% Bagian: turing-machines
% Bab: undecidability
% Subbagian: representing-tms

\documentclass[../../../include/open-logic-section]{subfiles}

\begin{document}

\olfileid[id]{tur}{und}{rep}
\olsection{Merepresentasikan Mesin Turing}

\begin{explain}
Untuk merepresentasikan mesin Turing dan perilakunya dengan !!a{sentence}
logika orde pertama, kita harus mendefinisikan bahasa yang sesuai. Bahasa itu
terdiri atas dua bagian: !!{predicate} untuk mendeskripsikan konfigurasi mesin,
serta ungkapan untuk memberi nomor pada langkah eksekusi (``saat'') dan posisi
pada pita.

Kita memperkenalkan dua jenis !!{predicate}, keduanya beraritas~2: untuk setiap
keadaan~$q$, !!a{predicate}~$\Obj Q_q$, dan untuk setiap simbol pita~$\sigma$,
!!a{predicate}~$\Obj S_\sigma$. Yang pertama memungkinkan kita mendeskripsikan
keadaan~$M$ dan posisi kepala pitanya; yang kedua memungkinkan kita
mendeskripsikan isi pita.

Untuk menyatakan posisi kepala pita dan banyaknya langkah yang telah
dijalankan, kita memerlukan cara untuk menyatakan bilangan. Hal ini dilakukan
dengan !!a{constant}~$\Obj 0$ dan fungsi beraritas~$1$, yaitu fungsi
penerus~$\prime$. Berdasarkan konvensi, fungsi itu dituliskan \emph{setelah}
argumennya (dan tanda kurungnya kita hilangkan).

Untuk setiap bilangan $n$ terdapat suku kanonik~$\num{n}$, yaitu
\emph{numeral} bagi~$n$, yang merepresentasikannya dalam~$\Lang L_M$.
$\num{0}$ adalah $\Obj 0$, $\num{1}$ adalah $\Obj 0'$, $\num{2}$ adalah
$\Obj 0''$, dan seterusnya. Secara lebih formal:
\begin{align*}
\num{0} & = \Obj 0 \\
\num{n+1} &= \num{n}'
\end{align*}
Suku $\num{0}$, yakni $\Obj 0$, menamai posisi paling kiri pada pita sekaligus
saat sebelum langkah eksekusi pertama (konfigurasi awal). Suku $\num{1}$,
yakni $\Obj 0'$, menamai petak di sebelah kanan petak paling kiri sekaligus
saat setelah langkah eksekusi pertama, dan seterusnya.

Kita juga memperkenalkan !!a{predicate}~$<$ untuk menyatakan baik urutan posisi
pita (ketika artinya ``di sebelah kiri'') maupun urutan langkah eksekusi
(ketika artinya ``sebelum'').

Setelah bahasanya tersedia, kita mendaftarkan ``aksioma-aksioma'' $!T(M,w)$,
yakni !!{sentence} yang secara bersama-sama mendeskripsikan perilaku~$M$ ketika
dijalankan pada masukan~$w$. Akan ada !!{sentence} yang menetapkan syarat pada
$\Obj 0$, $\prime$, dan $<$; !!{sentence} yang mendeskripsikan konfigurasi
masukan; serta !!{sentence} yang mendeskripsikan konfigurasi $M$ setelah mesin
menjalankan instruksi tertentu.
\end{explain}

\begin{defn}
  \ollabel{defn:tm-descr}
Diberikan mesin Turing $M = \tuple{Q, \Sigma, q_0, \delta}$, bahasa~$\Lang L_M$
terdiri atas:
\begin{enumerate}
\item !!^a{predicate} beraritas~2, $\Obj Q_q(x, y)$, untuk setiap keadaan
  $q \in Q$. Secara intuitif, $\Obj Q_q(\num{m}, \num{n})$ menyatakan
  ``setelah $n$ langkah, $M$ berada dalam keadaan~$q$ dan memindai petak
  ke-$m$.''
\item !!^a{predicate} beraritas~2, $\Obj S_\sigma(x, y)$, untuk setiap simbol
  $\sigma\in \Sigma$. Secara intuitif, $\Obj S_\sigma(\num{m}, \num{n})$
  menyatakan ``setelah $n$ langkah, petak ke-$m$ memuat simbol~$\sigma$.''
\item !!^a{constant} $\Obj 0$
\item !!^a{function} beraritas~1, $\prime$
\item !!^a{predicate} beraritas~2, $<$
\end{enumerate}
\end{defn}

!!^{sentence}s yang mendeskripsikan operasi mesin Turing~$M$ pada masukan
$w = \sigma_{i_1}\dots\sigma_{i_k}$ adalah sebagai berikut:
\begin{enumerate}
\item Aksioma yang mendeskripsikan bilangan dan~$<$:
\begin{enumerate}
\item !!^a{sentence} yang menyatakan bahwa setiap bilangan lebih kecil daripada
penerusnya:
\[
\lforall[x][x < x']
\]
\item !!^a{sentence} yang memastikan bahwa $<$ bersifat transitif:
\[
\lforall[x][\lforall[y][\lforall[z][
      ((x < y \land y < z) \lif x < z)]]]
\]
\end{enumerate}
\item Aksioma yang mendeskripsikan konfigurasi masukan:
\begin{enumerate}
\item Setelah $0$ langkah---sebelum mesin mulai---$M$ berada dalam keadaan
  awal~$q_0$ dan memindai petak~$1$:
\[
\Obj Q_{q_0}(\num{1}, \num{0})
\]
\item $k+1$ petak pertama memuat simbol $\TMendtape$,
  $\sigma_{i_1}$, \dots, $\sigma_{i_k}$:
\[
\Obj S_\TMendtape(\num{0}, \num{0}) \land
\Obj S_{\sigma_{i_1}}(\num{1}, \num{0}) \land
\dots \land
\Obj S_{\sigma_{i_k}}(\num{k}, \num{0})
\]
\item Selain itu, pita kosong:
\[
\lforall[x][(\num{k} < x \lif \Obj S_\TMblank(x, \num{0}))]
\]
\end{enumerate}
\item Aksioma yang mendeskripsikan transisi dari satu konfigurasi ke
  konfigurasi berikutnya:

Untuk uraian berikut, misalkan $!A(x, y)$ adalah konjungsi semua !!{formula}
berbentuk
\[
\lforall[z][
  (((z < x \lor x < z) \land \Obj S_\sigma(z, y))
  \lif \Obj S_\sigma(z, y'))]
\]
dengan $\sigma \in \Sigma$. Kita menggunakan $!A(\num{m},\num{n})$ untuk
menyatakan ``selain pada petak~$m$, pita setelah $n+1$ langkah sama dengan pita
setelah $n$ langkah.''
\begin{enumerate}
\item \ollabel{rep-right} Untuk setiap instruksi $\delta(q_i, \sigma) =
  \tuple{q_j, \sigma', \TMright}$, !!{sentence} tersebut adalah:
\begin{align*}
& \lforall[x][\lforall[y][(
   (\Obj Q_{q_i}(x, y) \land \Obj S_{\sigma}(x, y)) \lif {}]] \\
&\qquad   (\Obj Q_{q_j}(x', y') \land \Obj S_{\sigma'}(x, y') \land
!A(x, y)))
\end{align*}
Artinya, jika setelah~$y$ langkah mesin berada dalam keadaan~$q_i$ dan
memindai petak~$x$ yang memuat simbol~$\sigma$, maka setelah $y+1$ langkah
mesin memindai petak~$x+1$, berada dalam keadaan~$q_j$, petak~$x$ kini memuat
$\sigma'$, dan setiap petak selain~$x$ memuat simbol yang sama seperti setelah
$y$ langkah.

\item \ollabel{rep-left} Untuk setiap instruksi $\delta(q_i, \sigma) =
  \tuple{q_j, \sigma', \TMleft}$, !!{sentence} tersebut adalah:
\begin{align*}
& \lforall[x][\lforall[y][
    ((\Obj Q_{q_i}(x', y) \land \Obj S_{\sigma}(x', y)) \lif {}]]\\
& \qquad   (\Obj Q_{q_j}(x, y') \land \Obj S_{\sigma'}(x', y') \land
!A(x', y))) \land {}\\
& \lforall[y][((\Obj Q_{q_i}(\num{0}, y) \land \Obj S_{\sigma}(\num{0},
    y)) \lif {}]\\
& \qquad (\Obj Q_{q_j}(\num{0}, y') \land \Obj S_{\sigma'}(\num{0},
  y') \land !A(\num{0}, y)))
\end{align*}
Luangkan waktu sejenak untuk memahami cara kerjanya. Sekarang kita tidak mulai
dengan ``jika mesin memindai petak~$x$ \dots'', melainkan ``jika mesin memindai
petak $x+1$ \dots''. Gerakan ke kiri berarti bahwa pada langkah berikutnya
mesin memindai petak~$x$. Akan tetapi, petak yang ditulisi adalah~$x+1$.
Kita melakukannya dengan cara ini karena kita tidak memiliki pengurangan atau
fungsi pendahulu.

Perhatikan bahwa bilangan berbentuk $x+1$ adalah $1$, $2$, \dots. Jadi, hal ini
tidak mencakup kasus ketika mesin memindai petak~$0$ dan harus bergerak ke kiri
(yang tentu saja tidak dapat dilakukannya---mesin hanya diam di tempat).
Kasus khusus itu dicakup oleh konjungsi kedua: jika setelah $y$ langkah mesin
memindai petak~$0$ dalam keadaan $q_i$ dan petak~$0$ memuat simbol~$\sigma$,
maka setelah $y+1$ langkah mesin masih memindai petak~$0$, kini berada dalam
keadaan~$q_j$, simbol pada petak~$0$ adalah $\sigma'$, dan petak selain
petak~$0$ memuat simbol yang sama seperti setelah $y$ langkah.
\item \ollabel{rep-stay} Untuk setiap instruksi $\delta(q_i, \sigma) =
  \tuple{q_j, \sigma', \TMstay}$, !!{sentence} tersebut adalah:
\begin{align*}
& \lforall[x][\lforall[y][(
   (\Obj Q_{q_i}(x, y) \land \Obj S_{\sigma}(x, y)) \lif {}]] \\
&\qquad   (\Obj Q_{q_j}(x, y') \land \Obj S_{\sigma'}(x, y') \land
!A(x, y)))
\end{align*}
\end{enumerate}
\end{enumerate}
Misalkan $!T(M, w)$ adalah konjungsi semua !!{sentence} di atas bagi mesin
Turing~$M$ dan masukan~$w$.

Untuk menyatakan bahwa~$M$ pada akhirnya berhenti, kita harus menemukan
!!a{sentence} yang menyatakan ``setelah sejumlah langkah, fungsi transisi tidak
terdefinisi.'' Misalkan $X$ adalah himpunan semua pasangan
$\tuple{q, \sigma}$ sedemikian sehingga~$\delta(q, \sigma)$ tidak terdefinisi.
Misalkan $!E(M, w)$ adalah !!{sentence}
\[
\lexists[x][\lexists[y][(\bigvee_{\tuple{q, \sigma} \in
      X}(\Obj Q_q(x, y) \land \Obj S_\sigma(x, y)))]]
\]

Jika kita menggunakan mesin Turing dengan keadaan penghentian khusus~$h$,
caranya bahkan lebih mudah: !!{sentence}~$!E(M, w)$
\[
\lexists[x][\lexists[y][\Obj Q_h(x, y)]]
\]
menyatakan bahwa mesin pada akhirnya berhenti.

\begin{prop}
\ollabel{prop:mlessk}
Jika $m < k$, maka $!T(M, w) \Entails \num{m} < \num{k}$.
\end{prop}

\begin{proof}
Latihan.
\end{proof}

\begin{prob}
Buktikan \olref[tur][und][rep]{prop:mlessk}.
(Petunjuk: gunakan induksi pada $k-m$.)
\end{prob}

\end{document}
